La densite de puissance liee a un operateur correspondant a l'action
d'une processus $\mathcal{O}$ dans le plasma au temps $t$ etant
donnee par \[
P^{(t)\mathcal{O}}=2\pi\int_{-1}^{1}d\xi\int_{0}^{\infty}dp\;\frac{p^{3}}{\gamma m_{e}}S_{p}^{(t)\mathcal{O}}\]
 celle-ci vaut sur la grille des flux en $p$

\[
P^{(k)\mathcal{O}}=2\pi\sum_{i=0}^{n_{p}-1}\;\sum_{j=0}^{n_{\xi}-1}\frac{p_{i}^{3}}{\gamma_{i}m_{e}}S_{p,i,j+1/2}^{(k)\mathcal{O}}\Delta p_{i}\Delta\xi_{j+1/2}\]
 ou le flux dans l'espace des impulsions est la somme de trois termes
\[
\left.S_{p}^{(k)}\right|_{i,j+1/2}^{\mathcal{O}}=-D_{pp,i,j+1/2}^{\mathcal{O}}\left.\frac{\partial f_{0}^{(k)}}{\partial p}\right|_{i,j+1/2}+D_{p\xi,i,j+1/2}^{\mathcal{O}}\frac{\sqrt{1-\xi_{j+1/2}^{2}}}{p_{i}}\left.\frac{\partial f_{0}^{(k)}}{\partial\xi}\right|_{i,j+1/2}+F_{p,i,j+1/2}^{\mathcal{O}}f_{0,i,j+1/2}^{(k)}\]


La discretisation complete des flux intervient pour le calcul de la
puissance. \begin{eqnarray*}
\left.S_{p}\right|_{i+1,j+1/2}^{(k)} & = & -D_{pp,i+1,j+1/2}\left.\frac{\partial f_{0}}{\partial p}\right|_{i+1,j+1/2}^{(k)}\\
 &  & +D_{p\xi,i+1,j+1/2}\frac{\sqrt{1-\xi_{j+1/2}^{2}}}{p_{i+1}}\left.\frac{\partial f_{0}}{\partial\xi}\right|_{i+1,j+1/2}^{(k)}+F_{p,i+1,j+1/2}f_{0,i+1,j+1/2}^{(k)}\\
\left.S_{p}\right|_{i,j+1/2}^{(k)} & = & -D_{pp,i,j+1/2}\left.\frac{\partial f_{0}}{\partial p}\right|_{i,j+1/2}^{(k)}\\
 &  & +D_{p\xi,i,j+1/2}\frac{\sqrt{1-\xi_{j+1/2}^{2}}}{p_{i}}\left.\frac{\partial f_{0}}{\partial\xi}\right|_{i,j+1/2}^{(k)}+F_{p,i,j+1/2}f_{0,i,j+1/2}^{(k)}\\
\left.S_{\xi}\right|_{i+1/2,j+1}^{(k)} & = & -D_{\xi p,i+1/2,j+1}\left.\frac{\partial f_{0}}{\partial p}\right|_{i+1/2,j+1}^{(k)}\\
 &  & +D_{\xi\xi,i+1/2,j+1}\frac{\sqrt{1-\xi_{j+1}^{2}}}{p_{i+1/2}}\left.\frac{\partial f_{0}}{\partial\xi}\right|_{i+1/2,j+1}^{(k)}+F_{\xi,i+1/2,j+1}f_{0,i+1/2,j+1}^{(k)}\\
\left.S_{\xi}\right|_{i+1/2,j}^{(k)} & = & -D_{\xi p,i+1/2,j}\left.\frac{\partial f_{0}}{\partial p}\right|_{i+1/2,j}^{(k)}\\
 &  & +D_{\xi\xi,i+1/2,j}\frac{\sqrt{1-\xi_{j}^{2}}}{p_{i+1/2}}\left.\frac{\partial f_{0}}{\partial\xi}\right|_{i+1/2,j}^{(k)}+F_{\xi,i+1/2,j}f_{0,i+1/2,j}^{(k)}\end{eqnarray*}


L'expression du flux peut etre mise sous la forme discrete suivante
a partir des expressions numeriques des derivees \begin{eqnarray*}
\left.S_{p}^{\left(k\right)}\right|_{i,j+1/2}^{\mathcal{O}} & = & -\frac{D_{pp,i,j+1/2}^{\mathcal{O}}}{\Delta p_{i}}f_{0,i+1/2,j+1/2}^{(k)}\\
 &  & +\frac{D_{pp,i,j+1/2}^{\mathcal{O}}}{\Delta p_{i}}f_{0,i-1/2,j+1/2}^{(k)}\\
 &  & +D_{p\xi,i,j+1/2}^{\mathcal{O}}\frac{\sqrt{1-\xi_{j+1/2}^{2}}}{p_{i}\left(\Delta\xi_{j+1}+\Delta\xi_{j}\right)}f_{0,i,j+3/2}^{(k)}\\
 &  & +D_{p\xi,i,j+1/2}^{\mathcal{O}}\frac{\sqrt{1-\xi_{j+1/2}^{2}}}{p_{i}\left(\Delta\xi_{j+1}+\Delta\xi_{j}\right)}f_{0,i,j-1/2}^{(k)}\\
 &  & +F_{p,i,j+1/2}^{\mathcal{O}}f_{0,i,j+1/2}^{(k)}\end{eqnarray*}
 sachant que $\left.\frac{\partial f_{0}}{\partial\xi}\right|_{i,j+1/2}^{(k)}=\frac{f_{0,i,j+3/2}^{(k)}-f_{0,i,j-1/2}^{(k)}}{\Delta\xi_{j+1}+\Delta\xi_{j}}$.
Sachant par ailleurs que \[
\left\{ \begin{array}{l}
f_{0,i,j+3/2}^{(k)}=\left(1-\delta_{p,i,j+3/2}\right)f_{0,i+1/2,j+3/2}^{(k)}+\delta_{p,i,j+3/2}f_{0,i-1/2,j+3/2}^{(k)}\\
f_{0,i,j+1/2}^{(k)}=\left(1-\delta_{p,i,j+1/2}\right)f_{0,i+1/2,j+1/2}^{(k)}+\delta_{p,i,j-1/2}f_{0,i-1/2,j+1/2}^{(k)}\\
f_{0,i,j-1/2}^{(k)}=\left(1-\delta_{p,i,j-1/2}\right)f_{0,i+1/2,j-1/2}^{(k)}+\delta_{p,i,j-1/2}f_{0,i-1/2,j-1/2}^{(k)}\end{array}\right.\]
 on obtient finalement

\begin{eqnarray*}
\left.S_{p}^{\left(k\right)}\right|_{i,j+1/2}^{\mathcal{O}} & = & -\frac{D_{pp,i,j+1/2}^{\mathcal{O}}}{\Delta p_{i}}f_{i+1/2,j+1/2}^{(k)}\\
 &  & +\frac{D_{pp,i,j+1/2}^{\mathcal{O}}}{\Delta p_{i}}f_{i-1/2,j+1/2}^{(k)}\\
 &  & +D_{p\xi,i,j+1/2}^{\mathcal{O}}\frac{\sqrt{1-\xi_{j+1/2}^{2}}}{p_{i}\left(\Delta\xi_{j+1}+\Delta\xi_{j}\right)}\left(1-\delta_{p,i,j+3/2}\right)f_{0,i+1/2,j+3/2}^{(k)}\\
 &  & +D_{p\xi,i,j+1/2}^{\mathcal{O}}\frac{\sqrt{1-\xi_{j+1/2}^{2}}}{p_{i}\left(\Delta\xi_{j+1}+\Delta\xi_{j}\right)}\delta_{p,i,j+3/2}f_{0,i-1/2,j+3/2}^{(k)}\\
 &  & +D_{p\xi,i,j+1/2}^{\mathcal{O}}\frac{\sqrt{1-\xi_{j+1/2}^{2}}}{p_{i}\left(\Delta\xi_{j+1}+\Delta\xi_{j}\right)}\left(1-\delta_{p,i,j-1/2}\right)f_{0,i+1/2,j-1/2}^{(k)}\\
 &  & +D_{p\xi,i,j+1/2}^{\mathcal{O}}\frac{\sqrt{1-\xi_{j+1/2}^{2}}}{p_{i}\left(\Delta\xi_{j+1}+\Delta\xi_{j}\right)}\delta_{p,i,j-1/2}f_{0,i-1/2,j-1/2}^{(k)}\\
 &  & +F_{p,i,j+1/2}^{\mathcal{O}}\left(1-\delta_{p,i,j+1/2}\right)f_{i+1/2,j+1/2}^{(k)}\\
 &  & +F_{p,i,j+1/2}^{\mathcal{O}}\delta_{p,i,j+1/2}f_{i-1/2,j+1/2}^{(k)}\end{eqnarray*}



\subsubsection{Puissance moyennee sur la surface de flux}

Pour un calcul de densite de puissance, avec des surfaces de flux
concentriques, la moyenne sur une surface de flux est donnee par \begin{eqnarray*}
\left\langle P^{\mathcal{O}}\right\rangle  & = & \frac{1}{A_{\Psi}}\int_{S_{\Psi}}P^{\mathcal{O}}\\
 & = & \frac{1}{4\pi^{2}rR_{0}}\int_{-\pi}^{\pi}rd\theta\int_{0}^{2\pi}Rd\varsigma P^{\mathcal{O}}\\
 & = & \int_{-\pi}^{\pi}\frac{d\theta}{2\pi}\frac{R}{R_{0}}P^{\mathcal{O}}\\
 & = & \int_{-\pi}^{\pi}\frac{d\theta}{2\pi}(1+\epsilon\cos\theta)P^{\mathcal{O}}\\
 & = & (1+\epsilon)\int_{-\pi}^{\pi}\frac{d\theta}{2\pi}\frac{1}{\Psi(\theta)}P^{\mathcal{O}}\end{eqnarray*}
 et donc \[
\left\langle P^{\mathcal{O}}\right\rangle =2\pi\int_{0}^{\infty}dp\;\frac{p^{3}}{\gamma m_{e}}(1+\epsilon)\int_{-\pi}^{\pi}\frac{d\theta}{2\pi}\frac{1}{\Psi}\int_{-1}^{1}d\xi S_{p}^{\mathcal{O}}\]
 On peut alors rajouter la somme sur les directions opposees pour
les particules piegees car

\begin{eqnarray*}
\int_{-1}^{1}d\xi\left[\frac{1}{2}\sum_{\sigma=\pm1}\right]_{T}S_{p}^{\mathcal{O}} & = & \int_{-1}^{-\xi_{T}}d\xi S_{p}^{\mathcal{O}}+\int_{\xi_{T}}^{1}d\xi S_{p}^{\mathcal{O}}+\frac{1}{2}\int_{-\xi_{T}}^{\xi_{T}}d\xi\sum_{\sigma=\pm1}S_{p}^{\mathcal{O}}\\
 & = & \int_{-1}^{-\xi_{T}}d\xi S_{p}^{\mathcal{O}}+\int_{\xi_{T}}^{1}d\xi S_{p}^{\mathcal{O}}+\frac{1}{2}\int_{-\xi_{T}}^{\xi_{T}}d\xi\left[S_{p}^{\mathcal{O}}(\xi)+S_{p}^{\mathcal{O}}(-\xi)\right]\\
 & = & \int_{-1}^{-\xi_{T}}d\xi S_{p}^{\mathcal{O}}+\int_{\xi_{T}}^{1}d\xi S_{p}^{\mathcal{O}}+\int_{-\xi_{T}}^{\xi_{T}}d\xi S_{p}^{\mathcal{O}}(\xi)\\
 & = & \int_{-1}^{1}d\xi S_{p}^{\mathcal{O}}\end{eqnarray*}


On a alors \begin{eqnarray*}
\left\langle P^{\mathcal{O}}\right\rangle  & = & 2\pi\int_{0}^{\infty}dp\;\frac{p^{3}}{\gamma m_{e}}(1+\epsilon)\int_{-\pi}^{\pi}\frac{d\theta}{2\pi}\frac{1}{\Psi}\int_{-1}^{1}d\xi_{0}\frac{\Psi\xi_{0}}{\xi}H\left[\left|\xi_{0}\right|-\sqrt{1-\frac{1}{\Psi}}\right]\left[\frac{1}{2}\sum_{\sigma=\pm1}\right]_{T}S_{p}^{\mathcal{O}}\\
 & = & 2\pi\int_{0}^{\infty}dp\;\frac{p^{3}}{\gamma m_{e}}\int_{-1}^{1}d\xi_{0}(1+\epsilon)\left[\frac{1}{2}\sum_{\sigma=\pm1}\right]_{T}\int_{-\pi}^{\pi}\frac{d\theta}{2\pi}\frac{\xi_{0}}{\xi}H\left[\left|\xi_{0}\right|-\sqrt{1-\frac{1}{\Psi}}\right]S_{p}^{\mathcal{O}}\\
 & = & 2\pi\int_{0}^{\infty}dp\;\frac{p^{3}}{\gamma m_{e}}\int_{-1}^{1}d\xi_{0}(1+\epsilon)\left[\frac{1}{2}\sum_{\sigma=\pm1}\right]_{T}\int_{-\theta_{c}}^{\theta_{c}}\frac{d\theta}{2\pi}\frac{\xi_{0}}{\xi}S_{p}^{\mathcal{O}}\\
 & = & 2\pi\int_{0}^{\infty}dp\;\frac{p^{3}}{\gamma m_{e}}\int_{-1}^{1}d\xi_{0}(1+\epsilon)\lambda\left(\xi_{0}\right)\left\{ S_{p}^{\mathcal{O}}\right\} \end{eqnarray*}
 ou $\left\{ S_{p}^{\mathcal{O}}\right\} $ $\equiv S_{p}^{\mathcal{O}(0)}$
est la moyenne sur la trajectoire du flux dans l'espace des impulsions,
terme qui a ete calcule precedemment pour resoudre l'equation de diffusion
et determiner $f_{0}^{(0)}.$de sorte que \[
\left\langle P^{\mathcal{O}}\right\rangle ^{(t)}=2\pi\int_{0}^{\infty}dp\;\frac{p^{3}}{\gamma m_{e}}\int_{-1}^{1}d\xi_{0}(1+\epsilon)\lambda\left(\xi_{0}\right)S_{p}^{\mathcal{O}(0)(t)}\]
 ou la puissance est evaluee au temps $t$.

Sous forme discrete \[
\bigskip\left\langle P^{\mathcal{O}}\right\rangle ^{(k)}=2\pi\sum_{i=0}^{n_{p}-1}\;\sum_{j=0}^{n_{\xi}-1}\frac{p_{i}^{3}}{\gamma_{i}m_{e}}\lambda\left(\xi_{0j+1/2}\right)S_{p,i,j+1/2}^{\mathcal{O}(0)(k)}\Delta p_{i}\Delta\xi_{0j+1/2}\]
 avec \[
\left.S_{p}^{(k)}\right|_{i,j+1/2}^{\mathcal{O}(0)}=-D_{pp,i,j+1/2}^{\mathcal{O}(0)}\left.\frac{\partial f_{0}^{(0)(k)}}{\partial p}\right|_{i,j+1/2}+D_{p\xi_{0},i,j+1/2}^{\mathcal{O}(0)}\frac{\sqrt{1-\xi_{0j+1/2}^{2}}}{p_{i}}\left.\frac{\partial f_{0}^{(0)(k)}}{\partial\xi_{0}}\right|_{i,j+1/2}+F_{p,i,j+1/2}^{\mathcal{O}(0)}f_{0,i,j+1/2}^{(0)(k)}\]
 et on obtient une formulation analogue du flux dans l'espace des
impulsions au cas local, soit \begin{eqnarray*}
\left.S_{p}^{\left(k\right)}\right|_{i,j+1/2}^{\mathcal{O}(0)} & = & -\frac{D_{pp,i,j+1/2}^{\mathcal{O}(0)}}{\Delta p_{i}}f_{i+1/2,j+1/2}^{(0)(k)}\\
 &  & +\frac{D_{pp,i,j+1/2}^{\mathcal{O}(0)}}{\Delta p_{i}}f_{i-1/2,j+1/2}^{(0)(k)}\\
 &  & +D_{p\xi,i,j+1/2}^{\mathcal{O}(0)}\frac{\sqrt{1-\xi_{0j+1/2}^{2}}}{p_{i}\left(\Delta\xi_{0j+1}+\Delta\xi_{0j}\right)}\left(1-\delta_{p,i,j+3/2}^{(0)}\right)f_{0,i+1/2,j+3/2}^{(0)(k)}\\
 &  & +D_{p\xi,i,j+1/2}^{\mathcal{O}(0)}\frac{\sqrt{1-\xi_{j+1/2}^{2}}}{p_{i}\left(\Delta\xi_{0j+1}+\Delta\xi_{0j}\right)}\delta_{p,i,j+3/2}^{(0)}f_{0,i-1/2,j+3/2}^{(0)(k)}\\
 &  & +D_{p\xi,i,j+1/2}^{\mathcal{O}(0)}\frac{\sqrt{1-\xi_{0j+1/2}^{2}}}{p_{i}\left(\Delta\xi_{0j+1}+\Delta\xi_{0j}\right)}\left(1-\delta_{p,i,j-1/2}^{(0)}\right)f_{0,i+1/2,j-1/2}^{(0)(k)}\\
 &  & +D_{p\xi,i,j+1/2}^{\mathcal{O}(0)}\frac{\sqrt{1-\xi_{0j+1/2}^{2}}}{p_{i}\left(\Delta\xi_{0j+1}+\Delta\xi_{0j}\right)}\delta_{p,i,j-1/2}f_{0,i-1/2,j-1/2}^{(0)(k)}\\
 &  & +F_{p,i,j+1/2}^{\mathcal{O}(0)}\left(1-\delta_{p,i,j+1/2}^{(0)}\right)f_{0,i+1/2,j+1/2}^{(0)(k)}\\
 &  & +F_{p,i,j+1/2}^{\mathcal{O}(0)}\delta_{p,i,j+1/2}^{(0)}f_{0,i-1/2,j+1/2}^{(0)(k)}\end{eqnarray*}


On peut ainsi calculer la contribution de chaque processus, collisions,
champ electrique inductif, ondes RF.
